%! AP Statistics — Unit 4, Lesson 4.8 — Lesson Plan
\documentclass[10pt]{article}
\usepackage{apstats-article}
\usepackage{apstats-boxes}
\usepackage{graphicx}
\graphicspath{{images/}}
\newcommand{\TallMath}[1]{$\displaystyle #1\rule[-1.4em]{0pt}{3.2em}$}
\newcommand{\UnitNumberName}{Unit 4: Probability, Random Variables, and Probability Distributions \quad}
\newcommand{\LessonNumberName}{Lesson 4.8: Mean and Standard Deviation of Random Variables}

\begin{document}
\begin{center}
    {\Large\bfseries \CourseName: \SchoolYear \\
    {\small\bfseries \UnitNumberName \LessonNumberName}}
\end{center}

\begin{tcolorbox}[colback=sky, colframe=navy, boxrule=0.9pt, arc=2mm,
    left=2mm, right=2mm, top=1.5mm, bottom=1.5mm]
    \textbf{Primary Objective:} Students will calculate and interpret the mean (expected value)
    and standard deviation of a discrete random variable, using the computation table and
    connecting both statistics to the long-run behavior of the variable in context
    (Big Idea: VAR-5).
\end{tcolorbox}

\renewcommand{\arraystretch}{1.6}
\begin{skillbox}[Priority Ideas \& Skills]{goldbox}
\begin{minipage}[t]{0.48\linewidth}
    \textbf{AP Skill 3 --- Using Probability and Simulation}
    \begin{itemize}
        \item Compute $\mu_X = \sum x \cdot P(X=x)$ (Skill 3.B; VAR-5.C).
        \item Compute $\sigma^2_X = \sum (x-\mu_X)^2 \cdot P(X=x)$ and
              $\sigma_X = \sqrt{\sigma^2_X}$ (Skill 3.B; VAR-5.C).
        \item Interpret $\mu_X$ and $\sigma_X$ in context (Skill 4.B; VAR-5.D).
    \end{itemize}
\end{minipage}
\hfill
\begin{minipage}[t]{0.48\linewidth}
    \textbf{Key Understandings}
    \begin{itemize}
        \item $\mu_X$ is the long-run average outcome if the random process is repeated
              many times (not necessarily a possible value of $X$).
        \item $\sigma_X$ measures how much the values of $X$ typically deviate from
              the mean.
        \item AP exams reward contextual interpretation: students must say what
              $\mu_X$ and $\sigma_X$ mean for the real-world situation.
        \item The computation table organizes the weighted sums systematically.
    \end{itemize}
\end{minipage}
\end{skillbox}

\begin{skillbox}[{Vocabulary, Concepts, \& Theorems}]{greenbox}
\begin{minipage}[t]{0.48\linewidth}
    \begin{itemize}
        \item \textbf{Expected value} ($\mu_X$) --- $\mu_X = \sum x \cdot P(X=x)$;
              the long-run average
        \item \textbf{Variance} ($\sigma^2_X$) --- $\sigma^2_X = \sum (x-\mu_X)^2 \cdot P(X=x)$
        \item \textbf{Standard deviation} ($\sigma_X$) --- $\sqrt{\sigma^2_X}$
    \end{itemize}
\end{minipage}
\hfill
\begin{minipage}[t]{0.48\linewidth}
    \begin{itemize}
        \item \textbf{Long-run average} --- the value that the sample mean approaches
              over many independent repetitions (Law of Large Numbers)
        \item \textbf{Computation table} --- systematic layout of columns
              $x$, $P(X=x)$, $x P(X=x)$, $(x-\mu_X)^2 P(X=x)$
    \end{itemize}
\end{minipage}
\end{skillbox}

\begin{skillbox}[Activate Prior Knowledge \& Spiral Review]{sky}
    \begin{minipage}[t]{0.58\linewidth}\vspace{0pt}
        \textbf{Prior Knowledge Check (3--4 min)}
        \begin{itemize}
            \item Review: reading a probability distribution table (Lesson 4.7).
            \item Connect to data analysis: sample mean $\bar{x}=\sum x_i/n$ (Unit 1).
                  Today we generalize to a \emph{weighted} average.
            \item Prompt: \textit{``A game pays \$0 or \$10 with equal probability.
                  What do you expect to win on average per play?''}
        \end{itemize}
        \textbf{Connections Forward}
        \begin{itemize}
            \item $\mu_{X+Y} = \mu_X + \mu_Y$ $\to$ Lesson 4.9
            \item Binomial mean $np$ and SD $\sqrt{np(1-p)}$ $\to$ Lesson 4.10
        \end{itemize}
    \end{minipage}
    \hfill
    \begin{minipage}[t]{0.38\linewidth}\vspace{0pt}\centering
        % Authored warm-up; see warmup/main.tex for spiral review problems.
    \end{minipage}
\end{skillbox}

\begin{skillbox}[Hook]{sky}
\begin{minipage}[t]{0.48\linewidth}
    \textbf{Hook (4--5 min)}\\
    A carnival game costs \$2 to play. Prizes:
    \begin{center}
    \small\renewcommand{\arraystretch}{1.3}
    \begin{tabular}{@{}ccc@{}}
    Prize ($x$) & \$0 & \$5 \\
    \hline
    Probability & 0.70 & 0.30 \\
    \end{tabular}
    \end{center}
    ``If you play 100 times, how much do you expect to win in total?
    Is playing worth the \$2 entry fee?''\\[4pt]
    Transition: \textit{``The expected value of the prize tells us the long-run
    payoff per play.''}
\end{minipage}
\hfill
\begin{minipage}[t]{0.48\linewidth}
    \textbf{Lesson Flow (15--18 min)}
    \begin{enumerate}
        \item Motivate $\mu_X$: a weighted average where the weights are probabilities.
        \item Build the computation table step by step.
        \item Compute $\mu_X$; interpret as a long-run average.
        \item Note: $\mu_X$ need not be a possible value of $X$.
        \item Add the variance column $(x-\mu_X)^2 \cdot P(X=x)$.
        \item Compute $\sigma^2_X$ and $\sigma_X$; interpret both.
        \item Emphasize: AP essays are graded on interpretation, not just arithmetic.
    \end{enumerate}
\end{minipage}
\end{skillbox}

\begin{skillbox}[Explicit Instruction: Computation Table]{sky}
\begin{minipage}[t]{0.48\linewidth}
    \textbf{Steps}
    \begin{enumerate}
        \item List every value $x$ and $P(X=x)$.
        \item Compute $x \cdot P(X=x)$; sum for $\mu_X$.
        \item For each row compute $(x - \mu_X)^2 \cdot P(X=x)$.
        \item Sum that column for $\sigma^2_X$.
        \item $\sigma_X = \sqrt{\sigma^2_X}$.
        \item Interpret both in context.
    \end{enumerate}
\end{minipage}
\hfill
\begin{minipage}[t]{0.48\linewidth}
    \textbf{Worked Example: Number of Cars per Household}
    \small
    \renewcommand{\arraystretch}{1.4}
    \begin{tabular}{@{}ccccc@{}}
    $x$ & $P$ & $xP$ & $(x{-}\mu)^2$ & $(x{-}\mu)^2 P$ \\
    \hline
    0 & .10 & .00 & 2.4025 & .2403 \\
    1 & .40 & .40 & 0.3025 & .1210 \\
    2 & .35 & .70 & 0.2025 & .0709 \\
    3 & .15 & .45 & 2.1025 & .3154 \\
    \end{tabular}
    $\mu_X = 1.55$; $\sigma^2_X = 0.7476$; $\sigma_X \approx 0.865$\\[4pt]
    \textit{Interpretation:} On average, a household owns 1.55 cars. The number of
    cars typically varies from this mean by about 0.87 cars.
\end{minipage}
\end{skillbox}

\begin{skillbox}[Active Monitoring]{redbox}
\begin{minipage}[t]{0.48\linewidth}
    \textbf{Circulate and Check}
    \begin{itemize}
        \item Watch for summing probabilities instead of $x \cdot P(X=x)$.
        \item Check that $(x - \mu_X)$ uses the unrounded $\mu_X$.
        \item Cold-call: ``$\mu_X = 1.55$ cars --- can a household own 1.55 cars?
              What does this number mean?''
        \item Cold-call: ``What does $\sigma_X \approx 0.87$ cars tell us?''
    \end{itemize}
\end{minipage}
\hfill
\begin{minipage}[t]{0.48\linewidth}
    \textbf{Common Errors}
    \begin{itemize}
        \item Forgetting to square $(x - \mu_X)$.
        \item Not taking the square root for $\sigma_X$.
        \item Interpretation that only restates the number without context.
        \item Rounding $\mu_X$ too early; carry extra decimal places.
    \end{itemize}
\end{minipage}
\end{skillbox}

\begin{skillbox}[Group Work \& Differentiation]{redbox}
\begin{minipage}[t]{0.48\linewidth}
    \textbf{Activity Summary (10 min)}\\
    Three-tier differentiated activity (see \texttt{activity/}):
    \begin{itemize}
        \item \textbf{Tier R:} Partially completed table; students compute $\mu_X$ only
              and write one interpretation sentence.
        \item \textbf{Tier A:} Full computation table; students compute $\mu_X$ and
              $\sigma_X$ and write two interpretations.
        \item \textbf{Tier E:} Two-game comparison using $\mu_X$ and $\sigma_X$;
              students justify which game to play.
    \end{itemize}
\end{minipage}
\hfill
\begin{minipage}[t]{0.48\linewidth}
    \textbf{Differentiation Notes}
    \begin{itemize}
        \item \textit{Support}: Provide the completed first row and give $\mu_X$ before
              the variance columns.
        \item \textit{ELL}: Anchor chart: ``Expected value = long-run average.''
        \item \textit{Extension}: Students create a ``fair game'' by adjusting a
              prize so that $\mu_X = 0$ (net winnings).
    \end{itemize}
\end{minipage}
\end{skillbox}

\begin{skillbox}[Individual Work \& Assessment]{redbox}
\begin{minipage}[t]{0.48\linewidth}
    \textbf{Exit Ticket (5 min)}\\
    $X$ = goals scored per game:
    \begin{center}
    \small\renewcommand{\arraystretch}{1.3}
    \begin{tabular}{@{}cccc@{}}
    $x$ & 0 & 1 & 2 \\
    \hline
    $P(X=x)$ & 0.30 & 0.50 & 0.20 \\
    \end{tabular}
    \end{center}
    \begin{enumerate}
        \item Compute $\mu_X$.
        \item Interpret $\mu_X$ in context.
        \item Compute $\sigma_X$.
    \end{enumerate}
\end{minipage}
\hfill
\begin{minipage}[t]{0.48\linewidth}
    \textbf{AP-Style Check}\\
    \textit{Which best interprets an expected value of $\mu_X = 2.3$ customers
    per hour?}
    \begin{enumerate}[label=(\Alph*)]
        \item Exactly 2.3 customers arrive each hour.
        \item The most likely number of customers per hour is 2.3.
        \item Over many hours, the average number of customers per hour will be
              close to 2.3.
        \item At least 2.3 customers arrive in any given hour.
    \end{enumerate}
    Answer: (C).
\end{minipage}
\end{skillbox}

\begin{skillbox}[Reinforcement \& Extension]{goldbox}
\begin{minipage}[t]{0.48\linewidth}
    \textbf{Homework / Practice}
    \begin{itemize}
        \item Five problems: compute $\mu_X$ and $\sigma_X$ from tables, interpret
              both statistics, and use expected value to make a decision.
        \item Strong emphasis on writing complete interpretation sentences.
    \end{itemize}
\end{minipage}
\hfill
\begin{minipage}[t]{0.48\linewidth}
    \textbf{Extension / Preview}
    \begin{itemize}
        \item \textit{Extension}: Compare two games; decide which is better given
              $\mu_X$ and $\sigma_X$ (risk vs.\ reward).
        \item \textit{Preview}: Lesson 4.9 --- when two independent random variables
              are combined, how do means and variances add?
    \end{itemize}
\end{minipage}
\end{skillbox}

\end{document}
